\cotefile{72}
\scanpage{35}
\margin{?}
D'ailleurs, si $N=2$, $\exists!$ \deleted{\unsure{quadrique}}\inserted{conique} $Q\subset X$ \unsure{invariante} par $\varphi^+$, et elle est $=$ invariante par $\varphi$.

[Peut-être faut-il se donner au préalable \deleted{\unsure{formes quadratiques}}\inserted{hyp. qu. entière lisse} invariante \unsure{...}]
\mspage{34}

